\documentclass[12pt]{article}
\usepackage{amsmath,amssymb,amsthm}
\usepackage[margin=1in]{geometry}
\usepackage{booktabs}

\newtheorem{theorem}{Theorem}
\newtheorem{lemma}[theorem]{Lemma}
\newtheorem{corollary}[theorem]{Corollary}
\theoremstyle{definition}
\newtheorem{definition}[theorem]{Definition}

\title{Problem 162: Hexadecimal Numbers}
\author{Project Euler Solution}
\date{}

\begin{document}
\maketitle

\section{Problem Statement}
In the hexadecimal number system (digits $0$--$9$ and $A$--$F$), how many hexadecimal numbers containing at most 16 digits exist with all of the digits $0$, $1$, and $A$ present at least once? Give the answer as an uppercase hexadecimal number.

\section{Definitions}

\begin{definition}
An \emph{$n$-digit hexadecimal number} is a string $d_1 d_2 \cdots d_n$ over $\{0,\ldots,9,A,\ldots,F\}$ with $d_1 \neq 0$.
\end{definition}

\section{Mathematical Development}

\begin{theorem}[Inclusion-exclusion formula]
Let $N(n)$ count $n$-digit hex numbers whose digit set contains $\{0,1,A\}$. Then:
\[
N(n) = 15 \cdot 16^{n-1} - 15^n - 2 \cdot 14 \cdot 15^{n-1} + 2 \cdot 14^n + 13 \cdot 14^{n-1} - 13^n.
\]
\end{theorem}

\begin{proof}
Define $A_0, A_1, A_A$ as the sets of $n$-digit hex numbers missing digit $0$, $1$, $A$ respectively. By inclusion-exclusion, $N(n) = |S_n| - |A_0 \cup A_1 \cup A_A|$ where $|S_n| = 15 \cdot 16^{n-1}$.

The individual and intersection counts are:
\begin{center}
\begin{tabular}{@{}lll@{}}
\toprule
Set & Digit pool (first / rest) & Count \\
\midrule
$A_0$ & $\{1,\ldots,F\}$ / same & $15^n$ \\
$A_1$ & $\{2,\ldots,F\} \setminus\{1\}$ / $\{0,\ldots,F\}\setminus\{1\}$ & $14 \cdot 15^{n-1}$ \\
$A_A$ & $\{1,\ldots,F\}\setminus\{A\}$ / $\{0,\ldots,F\}\setminus\{A\}$ & $14 \cdot 15^{n-1}$ \\
$A_0 \cap A_1$ & $\{2,\ldots,F\}$ / same & $14^n$ \\
$A_0 \cap A_A$ & $\{1,\ldots,9,B,\ldots,F\}$ / same & $14^n$ \\
$A_1 \cap A_A$ & $\{2,\ldots,9,B,\ldots,F\}$ / $+\{0\}$ & $13 \cdot 14^{n-1}$ \\
$A_0 \cap A_1 \cap A_A$ & $\{2,\ldots,9,B,\ldots,F\}$ / same & $13^n$ \\
\bottomrule
\end{tabular}
\end{center}
Substituting into the inclusion-exclusion formula yields the claimed expression.
\end{proof}

\begin{theorem}[Summation]
The answer is $T = \sum_{n=1}^{16} N(n)$.
\end{theorem}

\begin{proof}
We sum over all valid digit lengths. Note $N(1) = N(2) = 0$ since at least 3 digits are needed to contain three distinct values with a non-zero leading digit constraint.
\end{proof}

\begin{corollary}[Closed-form evaluation]
Each of the six terms in $N(n)$ yields a geometric series over $n = 1, \ldots, 16$, e.g., $\sum_{n=1}^{16} 15 \cdot 16^{n-1} = 16^{16} - 1$. All values fit in 64-bit unsigned integers.
\end{corollary}

\section{Complexity}
\begin{itemize}
    \item \textbf{Time:} $O(16)$ multiplications, or $O(1)$ via closed-form geometric series.
    \item \textbf{Space:} $O(1)$.
\end{itemize}

\section{Answer}

\[
\boxed{\texttt{3D58725572C62302}}
\]

\end{document}
